\section{Conclusion}\label{sec:con}

We formulated \emph{online harness learning} as a setting for studying how frozen LLM agents can improve from sequential, one-shot executions. Through \algname{}, we instantiated this setting with context-to-harness skill compilation, which turns noisy execution contexts into a reusable external harness for future task solving. Beyond proposing a pipeline, our goal is to use the harness as an analytical lens for understanding online self-improvement. This perspective moves self-evolving agent research beyond aggregate success rates toward a more principled study of how agents produce, organize, transfer, and apply reusable experience in realistic task streams.

\section*{Limitations}

Our study focuses on single-agent LLM systems in text-based, tool-use, web, command-line, and software-engineering environments. We do not evaluate embodied agents or multi-agent systems, where the harness may involve additional forms of interaction, coordination, and environmental feedback. We also instantiate the harness with natural-language guidance; other harness formats, such as executable code skills or structured programs, are not covered in this study.